\subsection{膜的 Floer 上同调}\label{sec:12e}
先引入一个初等记号。固定系数域 $\K$。对任意一维实向量空间 $\delta$，以 $|\delta|_{\K}$ 表示由 $\delta$ 的两个定向生成、并施加两个生成元之和为零这一关系的一维 $\K$ 向量空间。显然，$\R$ 线性同构 $c:\delta_1\to\delta_2$ 诱导 $\K$ 线性同构 $|c|_{\K}:|\delta_1|_{\K}\to|\delta_2|_{\K}$。分别称其为 $\delta$ 与 $c$ 的 $\K$ \emph{规范化}。规范化也适用于族，将实线丛变成纤维为 $\K$ 的局部系数系统，对它们之间的同构亦然。可将它理解为通过同态 $\R^\times\to\{\pm1\}\subset\K^\times$、$t\mapsto t/|t|$ 取关联丛。

升级后的 Floer 上同调为 $M$ 中任意一对正合 Lagrangian 膜 $(L_0^\#,\allowbreak L_1^\#)$ 关联有限维分次 $\K$ 向量空间 $HF^*(L_0^\#,\allowbreak L_1^\#)$；取 $\operatorname{char}(\K)=2$ 并将所有
\SourcePage{188}
分次部分直和，就还原出 $HF^{\mathrm{pr}}(L_0,\allowbreak L_1)$。下面回顾\sr{sec:8}{第~8~节}讨论的若干基本性质，并说明有膜结构时如何适配。
\begin{itemize}
\item \emph{三角形乘积。}存在分次、结合且含幺的乘积
\begin{equation}\label{eq:12.11}
HF^*(L_1^\#,L_2^\#)\otimes HF^*(L_0^\#,L_1^\#)\longrightarrow HF^*(L_0^\#,L_2^\#).
\end{equation}
\item \emph{改变膜结构。}若对所有涉及的膜施加同一个运算，Floer 上同调群不变，\eref{12.11}也不变。这里“运算”可以是给分次加上同一 $\Z$ 中常数，也可以是将各 Pin 结构与某个共同线丛 $\beta\to M$ 的限制作张量积。另有不那么显然的典范同构，它也说明“平移”这一名称合理：
\begin{equation}\label{eq:12.12}
HF^*(L_0^\#,S^\sigma L_1^\#)\cong HF^{*+\sigma}(L_0^\#,L_1^\#).
\end{equation}
它们也与乘积相容，即下图交换：
\SourceEquationDiagram{diagram-12-13}{0.55453786}{12.13}
\item \emph{PSS 同构。}存在典范同构 $HF^*(L^\#,\allowbreak L^\#)\cong H^*(L;\K)$，在 $L_0^\#=L_1^\#=L_2^\#=L^\#$ 时将\eref{12.11}化为通常的杯积。更一般地，若用任意两个线丛 $\beta_0,\allowbreak \beta_1$ 扭曲 $L^\#$，结果是 $L$ 的、取值于相应局部系数系统的上同调：
\begin{equation}\label{eq:12.14}
HF^*(L^\#\otimes\beta_0,L^\#\otimes\beta_1)\cong H^*(L;|\beta_0^\vee\otimes\beta_1|_{\K}).
\end{equation}
\item \emph{对偶性。}对任意 $L^\#$，Poincaré 对偶与 PSS 同构给出映射
\begin{equation}\label{eq:12.15}
\begin{aligned}
\int_{L^\#}:HF^n(L^\#,L^\#\otimes\lambda^{\mathrm{top}}(TL))
&\cong H^n(L;|\lambda^{\mathrm{top}}TL|_{\K})\\
&\cong H_0(L;\K)\longrightarrow\K.
\end{aligned}
\end{equation}
由此构造的次数为 $-n$ 的配对
\[
\begin{gathered}
HF^*(L_1^\#,L_0^\#\otimes\lambda^{\mathrm{top}}(TL_0))\otimes HF^*(L_0^\#,L_1^\#)\longrightarrow\K,\\
\langle a_2,a_1\rangle=\int_{L_0^\#}a_2a_1
\end{gathered}
\]
非退化且分次对称。后一个性质需要略作解释。利用选定同构 $\lambda^{\mathrm{top}}(TL_k)\cong\xi_M|L_k$，可将 $a_2$ 视为 $HF^*(L_1^\#,\allowbreak L_0^\#\otimes\xi_M|L_0)$ 的元素，将 $a_1$ 视为
\SourcePage{189}
$HF^*(L_0^\#\otimes\xi_M|L_0,\allowbreak L_1^\#\otimes\lambda^{\mathrm{top}}(TL_1))$ 的元素。于是 $a_1a_2\in HF^*(L_1^\#,\allowbreak L_1^\#\otimes\lambda^{\mathrm{top}}(TL_1))$，因此下式有意义：
\[
\langle a_1,a_2\rangle=(-1)^{|a_1|\cdot|a_2|}\langle a_2,a_1\rangle.
\]
\item \emph{同痕不变性。}从 $L_0^\#$ 到 $L_1^\#$ 的正合 Lagrangian 膜同痕给出 $HF^0(L_0^\#,\allowbreak L_1^\#)$ 的一个选定元素。这些元素的性质仍如\secref{8k}；特别地，可用来证明 $HF^*$ 具有同痕不变性。
\end{itemize}
\begin{definition*}
以正合 Lagrangian 膜为对象、Floer 上同调群为态射，并用\eref{12.11}给出复合，定义 \emph{Donaldson–Fukaya 范畴} $H\F(M)$。
\end{definition*}
它是 $\K$ 线性的 $\Z$ 分次范畴；取 $\operatorname{char}(\K)=2$ 并忘掉分次，便退化为 $H\F(M)^{\mathrm{pr}}$。由于 Poincaré 对偶配对具有扭曲形式，一般情况下它不是 Frobenius 范畴，除非 $\xi_M$ 平凡。\znote{180}{配对右端落在 $\lambda^{\mathrm{top}}(TL_0)$ 的规范化里，而不是平凡地落在 $\K$：Frobenius 范畴要求存在与复合相容的非退化配对 $\operatorname{hom}(A,B)\cong\operatorname{hom}(B,A)^\vee$，这里多出的扭转使该同构带上一个线丛因子，故只有 $\xi_M$ 平凡、扭转可平凡化时才成立。第 12j 节把这一缺陷代数化：Donaldson–Fukaya 范畴满足的是弱 Calabi–Yau 性质，对偶由 Serre 函子而非恒等函子实现，次数为 $-\dim M$。换言之，Frobenius 性质并未完全丢失，而是被 $\lambda^{\mathrm{top}}$ 的扭转吸收掉了。[来源：nLab “Serre functor”，https:/\allowbreak /\allowbreak ncatlab.org/\allowbreak nlab/\allowbreak show/\allowbreak Serre+\allowbreak functor（2026-11-15 访问）；并见原书 §12j。]}

\subsection{确定分次与符号}\label{sec:12f}
有了全部准备，$HF^*(L_0^\#,\allowbreak L_1^\#)$ 的定义相当直接。选择正则 Floer 数据，对每个 $k\in\Z$ 定义分次 Floer 上链群
\begin{equation}\label{eq:12.16}
CF^k(L_0^\#,L_1^\#)=\bigoplus_{i(y)=k}|o(y)|_{\K}.
\end{equation}
设 $u\in\M_Z(y_0,\allowbreak y_1)$ 为 Floer 方程的非驻定解。有显然的短正合列
\[
0\longrightarrow\R\longrightarrow T\M_Z(y_0,y_1)_u\longrightarrow T\M_Z^*(y_0,y_1)_{[u]}\longrightarrow0,
\]
其中 $\R$ 项对应平移方向；更精确地，约定其正生成元为 $\partial u/\partial s$。由此可识别 $\M_Z$ 与 $\M_Z^*$ 的切空间的最高次外幂，再由\eref{12.2}，得到 $\M_Z^*(y_0,\allowbreak y_1)$ 相对于 $o(y_0)\otimes o(y_1)^\vee$ 的选定定向。特别地，若 $[u]$ 是 $\M_Z^*(y_0,\allowbreak y_1)$ 的孤立点，就得到同构
\begin{equation}\label{eq:12.17}
c_u:o(y_1)\longrightarrow o(y_0).
\end{equation}
将边界算子 $\partial$ 的 $(y_0,\allowbreak y_1)$ 系数定义为所有这种 $c_u$ 的 $\K$ 规范化之和：
\begin{equation}\label{eq:12.18}
\partial_{y_0,y_1}=\sum_u|c_u|_{\K}:|o(y_1)|_{\K}\longrightarrow|o(y_0)|_{\K}.
\end{equation}
这是一种较精巧的、不依赖基的说法，意思是每个 $u$ 依照 $c_u$ 的符号贡献 $\pm1$。再看\eref{12.2}即可知，$\partial$ 的次数为所需的 $1$。\znote{181}{$\partial$ 的每个矩阵元都是数出来的：一个孤立 Floer 轨线给出同构 $c_u:o(y_1)\to o(y_0)$，规范化后贡献 $\pm1$。这样写的好处是符号自动一致：模空间的定向由 $o(y_0)\otimes o(y_1)^\vee$ 给出，而 $\partial^2=0$ 归结为紧化边界上两个端点的定向相反，下文用粘合映射的微分 $Dg$ 逐项验证。同一套“计数 + 行列式线”模式在 12g 定义 $\mu^d$、13c 定义 $\tilde\mu^d$ 时反复出现，值得在此处就熟悉。这种先数点、再取规范化和的写法，使 $\mu^d$ 的符号完全由模空间定向决定，而不依赖任何临时约定。[来源：Seidel, Fukaya categories and Picard–Lefschetz theory, §12f 与 §9j。]}

为说明\eref{12.18}合理，取一对解 $u_1\in\M_Z(y_0,\allowbreak \tilde y)$、$u_2\in\M_Z(\tilde y,\allowbreak y_1)$。将 $u_2$ 向右平移足够远后，
\SourcePage{190}
可粘合二者，得到 $u\in\M_Z(y_0,\allowbreak y_1)$。用 Gromov 紧化\eref{9.28}来说，$[u]$ 是一维模空间 $\M_Z^*(y_0,\allowbreak y_1)$ 的点，靠近紧化的边界点 $([u_1],\allowbreak [u_2])$。由粘合映射 $g$ 的结构，可见其微分 $Dg_{(u_1,\allowbreak u_2)}$ 具有以下性质：$(\partial u_1/\partial s,\allowbreak \partial u_2/\partial s)$ 映为平移切向量 $\partial u/\partial s$，而 $(-\partial u_1/\partial s,\allowbreak \partial u_2/\partial s)$ 映为远离边界点 $(u_1,\allowbreak u_2)$ 的向量。沿相反方向 $(\partial u_1/\partial s,\allowbreak -\partial u_2/\partial s)$ 移动 $(u_1,\allowbreak u_2)$，对应于将两个部分拉开，也就是令粘合长度趋于无穷。结合\eref{12.4}可知，通过同构\eref{12.2}由 $c_{u_1}\circ c_{u_2}:o(y_1)\to o(y_0)$ 得到的 $T\M_Z^*(y_0,\allowbreak y_1)_{[u]}$ 的定向，始终指向远离 Gromov 边界的方向。因此，同一连通分支的两个边界点贡献相反的映射 $|c_{u_1}\circ c_{u_2}|_{\K}$，从而 $\partial^2=0$。如通常一样，定义 $HF^*(L_0^\#,\allowbreak L_1^\#)=H(CF^*(L_0^\#,\allowbreak L_1^\#),\allowbreak \partial)$。

现取膜标记为 $(L_0^\#,\allowbreak \ldots,\allowbreak L_d^\#)$ 的 $(d+1)$ 点圆盘 $S$，选择正则 Floer 数据与扰动数据。在每个孤立点 $u\in\M_S(y_0,\allowbreak \ldots,\allowbreak y_d)$ 处，相对定向\eref{12.2}给出同构 $c_u:o(y_d)\otimes\cdots\otimes o(y_1)\to o(y_0)$。如\eref{12.18}，用它们定义映射
\[
C\Phi_S:CF^*(L_{d-1}^\#,L_d^\#)\otimes\cdots\otimes CF^*(L_0^\#,L_1^\#)\longrightarrow CF^*(L_0^\#,L_d^\#).
\]
为证明 $C\Phi_S$ 为链映射，考察 Gromov 紧化的两类边界点：
\begin{itemize}
\item 第一类为 $(u_1,\allowbreak [u_2])\in\M_S(y_0,\allowbreak \ldots,\allowbreak y_n,\allowbreak \tilde y,\allowbreak y_{n+2},\allowbreak \ldots,\allowbreak y_d)\times\M_Z^*(\tilde y,\allowbreak y_{n+1})$。平移向量 $(0,\allowbreak \partial u_2/\partial s)$ 在 $Dg$ 下的像是 $\M_S(y_0,\allowbreak \ldots,\allowbreak y_d)$ 的一个远离边界的切向量。因此，由 $c_{u_1}\circ c_{u_2}$ 通过\eref{12.2}得到的局部定向，与标准向外定向相差 $(-1)$ 乘以\eref{12.4}中的 Koszul 符号，其中 $\operatorname{index}(D_{S_2,\allowbreak u_2})=1$。
\item 第二类为 $([u_1],\allowbreak u_2)\in\M_Z^*(y_0,\allowbreak \tilde y)\times\M_S(\tilde y,\allowbreak y_1,\allowbreak \ldots,\allowbreak y_d)$。$(\partial u_1/\partial s,\allowbreak 0)$ 的像指向边界，且 Koszul 符号平凡。
\end{itemize}
合起来得到
\begin{equation}\label{eq:12.19}
\sum_n(-1)^*C\Phi_S(y_d,\ldots,\partial(y_{n+1}),y_n,\ldots,y_1)
=\partial(C\Phi_S(y_d,\ldots,y_1)),
\end{equation}
其中 $*=i(y_{n+2})+\cdots+i(y_d)$。沿同样思路作类似但更简单的论证，可知复合律\eref{8.2}成立，且不带任何符号。如\secref{8b}，利用 $C\Phi_S$ 诱导的上同调映射 $\Phi_S$，定义 $HF^0(L^\#,\allowbreak L^\#)$ 中的恒等态射（$d=0$）、三角形乘积（$d=2$）及延拓映射（$d=1$）。

也值得提一下，如何利用 $S$ 上扰动数据的族 $(K^s,\allowbreak J^s)$、$s\in[0;1]$ 构造链同伦。考虑
\SourcePage{191}
相应的参数化模空间，用其中的孤立点定义映射
\[
h:CF^*(L_{d-1}^\#,L_d^\#)\otimes\cdots\otimes CF^*(L_0^\#,L_1^\#)\longrightarrow CF^*(L_0^\#,L_d^\#)[-1].
\]
模空间的局部定向来自\eref{12.8}，粘合方面使用\eref{12.10}，取 $\RG=[0;1]$。与先前论证相比，唯一差别是第二类情形中 Koszul 约定引入另一个符号 $(-1)^{\dim(\RG)}=-1$。结果是
\begin{equation}\label{eq:12.20}
\begin{aligned}
C\Phi_S^1-C\Phi_S^0={}&\sum_n(-1)^*h(y_d,\ldots,\partial(y_{n+1}),y_n,\ldots,y_1)\\
&+\partial(h(y_d,\ldots,y_1)),
\end{aligned}
\end{equation}
其中 $C\Phi_S^k$ 是由 $(K^k,\allowbreak J^k)$ 导出的链映射，$*$ 如\eref{12.19}。这已足以建立基本的良定义性质。

沿\secref{12c}的讨论，只要 $\hat S$ 的每个边界圆周都有奇数个入向无穷远点，就可毫无困难地将理论推广到其他曲面 $S$，由此得到适当意义下的拓扑量子场论。若要允许不具备这一性质的曲面，则须对膜施加适当的定向假设，或使用\eref{12.6}的技巧。例如，取 $S=\bar H$，其端为出向端，膜标记为 $L^\#$。\eref{12.7}说明，每个孤立元素 $u\in\M_S(y)$ 都给出同构
\[
c_u:o(\Lambda_{y,0}^\#,\Lambda_{y,1}^\#\otimes\lambda^{\mathrm{top}}(\Lambda_{y,1}))\longrightarrow\K.
\]
将它们规范化后合并，便得到次数为 $-n$ 的链映射 $CF^*(L^\#,\allowbreak L^\#\otimes\lambda^{\mathrm{top}}(TL))\to\K$。诱导的上同调映射恰为\eref{12.15}。
